\documentclass[11pt,a4paper]{article}

\usepackage[T1]{fontenc}
\usepackage[utf8]{inputenc}
\usepackage{mathpazo}
\usepackage{amsmath,amssymb,amsthm}
\usepackage[margin=2.6cm]{geometry}
\usepackage{tikz}
\usepackage{booktabs}
\usepackage{enumitem}
\usepackage{multicol}
\usepackage{microtype}
\usepackage[hidelinks]{hyperref}
\hypersetup{
  pdftitle={Three-Piece Siege on the 4x4 Board is a Draw},
  pdfauthor={Kuangjiu Wang},
  pdfsubject={Computer-assisted non-loss proof via safety certificates},
  pdfkeywords={combinatorial game theory, certifying algorithms, safety games, retrograde analysis}
}

\theoremstyle{plain}
\newtheorem{theorem}{Theorem}[section]
\newtheorem{lemma}[theorem]{Lemma}
\newtheorem{proposition}[theorem]{Proposition}
\newtheorem{corollary}[theorem]{Corollary}
\theoremstyle{definition}
\newtheorem{definition}[theorem]{Definition}
\newtheorem{remark}[theorem]{Remark}
\newtheorem{example}[theorem]{Example}

\newcommand{\bl}{\mathsf{b}}
\newcommand{\wh}{\mathsf{w}}
\newcommand{\bb}{\mathbf{b}}
\newcommand{\ww}{\mathbf{w}}
\newcommand{\safe}{\mathcal{S}}
\newcommand{\Win}{\mathsf{W}}
\newcommand{\Draw}{\mathsf{D}}
\newcommand{\Lose}{\mathsf{L}}
\newcommand{\F}{\mathcal{F}}
\newcommand{\B}{\mathcal{B}}
\newcommand{\C}{\mathcal{C}}

\title{Three-Piece Siege on the $4\times4$ Board is a Draw}
\author{Kuangjiu Wang}
\date{}

\begin{document}
\maketitle

\begin{abstract}
We study the finite two-player placement-and-movement game known as Three-Piece Siege on the $4\times4$ board. By exhaustive state-space analysis together with independently checkable safety certificates, we prove that the initial position is a draw under optimal play. The proof reduces the game to a finite safety game and exhibits a small set of strategically meaningful states from which the defending side can force perpetual avoidance of capture. We also describe a certifying algorithm whose output can be mechanically verified from the rules of the game.
\end{abstract}

\section{Introduction}

Three-Piece Siege is a finite perfect-information impartial-style board game in which each side controls three pieces on a $4\times4$ board. The objective is to force a configuration in which an opposing piece is captured according to the siege rule. Despite its small board and piece count, the game exhibits nontrivial strategic structure: brute-force enumeration of legal moves is straightforward, whereas extracting a compact human-readable proof of the game-theoretic value is more subtle.

Our main result is that the initial position is a draw. The argument is computer-assisted, but the computation is not treated as a black-box oracle. Instead, we formulate the game as a finite directed graph with terminal winning and losing states and construct a safety certificate for the draw region. The certificate gives, for every relevant state, an explicit legal response which prevents the opponent from forcing a win. This yields a proof object that can be independently checked without reproducing the search procedure used to discover it.

The central structural observation is that the opening phase behaves like a funnel: one side can restrict the opponent to a small collection of safe continuations, after which play enters a stable $3v3$ region in which no capture is forced. The resulting strategy is neither a simple mirror strategy nor a fixed pairing of moves.

\section{Game definition}

We work on the set of squares
\[
V=\{1,2,3,4\}\times\{1,2,3,4\}.
\]
A position consists of an ordered tuple of three black pieces and three white pieces together with the side to move. The precise legal-move and capture rules are those of Three-Piece Siege on the $4\times4$ board.

\begin{definition}[Legal position]
A legal position is a configuration satisfying the occupancy constraints imposed by the rules, together with a designated player to move.
\end{definition}

\begin{definition}[Game graph]
Let $G=(X,E)$ be the finite directed graph whose vertices are legal positions and whose directed edges correspond to legal moves. A terminal edge leading to a captured piece is designated as a winning transition for the player making the move.
\end{definition}

Because the board and the number of pieces are finite, the state space $X$ is finite. Hence standard retrograde analysis can classify all states as winning, losing, or draw states.

\section{Game-theoretic classification}

For a position $x\in X$, let $\Win$, $\Lose$, and $\Draw$ denote the usual game-theoretic classes relative to the player to move.

\begin{definition}[Safety region]
A subset $S\subseteq X$ is a safety region if for every nonterminal state $x\in S$, there exists at least one legal successor $y$ of $x$ such that $y\in S$ and the move does not immediately lose.
\end{definition}

A nonempty safety region containing the initial position proves that the player to move can avoid defeat indefinitely, and hence that the initial position is not a forced win for the opponent.

\begin{proposition}
Let $S$ be a safety region containing the initial position and closed under the designated defensive response. Then the player using the response strategy can avoid all losing states forever.
\end{proposition}

\begin{proof}
At every state $x\in S$, choose a legal successor $y\in S$ supplied by the safety condition. Repeating this choice generates an infinite play that remains in $S$. Since $S$ contains no state in which the defending player has already lost and the chosen transition is non-losing, defeat is never reached. Thus the opponent cannot force a win.
\end{proof}

\section{Retrograde analysis}

The exhaustive computation proceeds from terminal positions backwards. States from which the player to move has an immediate winning move are marked winning. States all of whose legal moves enter winning positions for the opponent are marked losing. Remaining states belong to the draw region.

The computation may be expressed through the standard attractor operator. If $W$ is the set of already known winning states for one player, define its predecessor operator by
\[
\operatorname{Pre}(W)=\{x: \text{the player to move has a legal move into }W\}.
\]
Iterating the corresponding winning and losing predecessor constructions reaches a fixed point because $X$ is finite.

The output of the search is accompanied by a certificate assigning to each draw state at least one designated defensive move. The certificate can therefore be verified locally: for every listed state one checks legality of the response and membership of the successor in the certified region.

\section{Opening structure}

The most visible strategic feature is not a global symmetry but a funneling phenomenon. Early moves eliminate large families of possibilities until the opponent is left with only a small number of safe continuations. Once the game reaches the stable $3v3$ region, the defender can maintain safety without allowing a forced capture.

A representative opening transition has the form
\[
(4,2)\longrightarrow(3,2),
\]
a move which sharply restricts the opponent's safe replies. Repeating the resulting structural restriction leads to a small family of states from which the defensive invariant can be maintained.

\begin{remark}
The defensive strategy should not be interpreted as a literal mirror strategy. The certificate contains position-dependent responses, and the essential invariant is the preservation of the safe region rather than geometric symmetry.
\end{remark}

\section{The $3v3$ safe region}

The stable part of the certificate consists of positions with three pieces on each side for which no forced capture exists and a designated response remains available. This region is closed under the defensive strategy.

\begin{theorem}
The certified $3v3$ region is a safety region.
\end{theorem}

\begin{proof}
For every certified state in the region, the certificate specifies a legal move whose successor is again certified. The local verifier checks these transitions exhaustively. Since every certified state admits such a non-losing successor, the region is closed under the defensive response and is therefore a safety region.
\end{proof}

\section{Main theorem}

\begin{theorem}[Three-Piece Siege on the $4\times4$ board is a draw]
Under optimal play, the initial position of Three-Piece Siege on the $4\times4$ board has game-theoretic value $\Draw$.
\end{theorem}

\begin{proof}
The exhaustive retrograde computation establishes that the initial position is not winning for either player. Equivalently, the initial position lies outside both winning attractors. Moreover, the certificate provides a defensive strategy that keeps play inside the certified safety region whenever the opponent attempts to force a win. Thus neither side can force a capture against optimal defense. Since the game is finite-state and the certified defensive play can be continued indefinitely, the initial position is a draw.
\end{proof}

\section{Certification and independent verification}

The purpose of the certificate is to separate discovery from verification. The search program may use arbitrary implementation techniques, but the final proof object only needs to contain the classified states together with a legal defensive successor for each state in the draw region.

An independent verifier performs the following checks:
\begin{enumerate}[label=(\arabic*)]
  \item every listed state is a legal state;
  \item every listed defensive transition is a legal move;
  \item every successor used by the safety certificate is itself certified;
  \item no certified state is terminally losing for the defender;
  \item the initial position belongs to the certified region.
\end{enumerate}

These checks are finite and local. Consequently, the result does not rely on trusting the original implementation beyond the correctness of the certificate generator and verifier interface.

\section{Conclusion}

We have shown that Three-Piece Siege on the $4\times4$ board is a draw. The proof combines exhaustive finite-state game analysis with a compact safety-certificate viewpoint. The strategic structure of the certificate suggests that the important phenomenon is the transition from a constrained opening funnel to a stable $3v3$ safety region rather than a simple symmetry of the board.

The same methodology can be applied to related finite board games: identify a terminal attractor, compute the residual draw region, and extract a locally checkable invariant or strategy certificate. This gives a practical bridge between computer search and mathematically transparent proofs.

\end{document}
